% GENERATED FILE. Edit canonical lessons, then run npm run generate:books.
% canonical-source-hash: fnv1a64:8a27b4668dcd6bbe
% canonical-lessons: MR-C319-rosnai, MR-C319-rangoli, MR-C319-toran, MR-C319-phataka, MR-C319-atasbaji, MR-R319-review-312-319

\chapter{Lights and Colours}
\label{ch:mr-lights-and-colours}
\hlchaptermodality{319}

\begin{chapteropening}
\textbf{By the end of this chapter:} \emph{I can describe festive lights, rangoli and fireworks.}

Complete the last lesson of chapter 319: I can describe festive lights, rangoli and fireworks.
\end{chapteropening}

\section[\texorpdfstring{\mr{रोषणाई} (roṣṇāī)}{रोषणाई (roṣṇāī)}]{\mr{रोषणाई} (roṣṇāī) — festive lights; illuminations}

\label{lesson:MR-C319-rosnai}



\begin{quote}
Before the new one: say the Marathi for a thread ceremony, then the Marathi for an invitation card.
\end{quote}

\subsection*{You'll want to know: \mr{रोषणाई}}
\textbf{\mr{रोषणाई}} — \emph{roṣṇāī} — \textquotedblleft{}festive lights; illuminations\textquotedblright{}.

\begin{quote}
\mr{सणाला} \mr{संपूर्ण} \mr{गल्लीत} \mr{रोषणाई} \mr{केली} \mr{होती}.
\end{quote}

\begin{quote}
\emph{saṇālā sampūrṇa gallīt roṣṇāī kelī hotī.}
\end{quote}

\begin{quote}
For the festival the whole lane had been lit up.
\end{quote}

Feminine.

\mr{रोषणाई} \mr{करणे}, to put up lights. The letter \mr{ऋ} opens \mr{ऋतू}, a season; every \mr{ऋतू} has its festival.

\begin{cousinweb}
\textbf{\mr{रोषणाई}} (\emph{roṣṇāī}) is a loan from Persian \emph{rošanā'ī}, 'brightness, light', built on \emph{rošan}, 'bright'.
\end{cousinweb}

\begin{grammarlens}[title={What you've built}]
One more word for telling what happened.
\end{grammarlens}

\subsection*{Your turn}
\begin{itemize}
\raggedright
  \item \emph{Say it:} \emph{roṣṇāī}
  \item \emph{Say it:} \emph{roṣṇāī}, once more
  \item \emph{Say it:} \emph{saṇālā sampūrṇa gallīt roṣṇāī kelī hotī.} — For the festival the whole lane had been lit up.
  \item \emph{Recall:} say the Marathi for a thread ceremony, then the Marathi for an invitation card, then say \emph{roṣṇāī} again
\end{itemize}

\begin{tcolorbox}[breakable,colback=teal!4,colframe=teal!35!black,title={Before you move on}]
What does \mr{रोषणाई} mean? (\textquotedblleft{}Festive lights; illuminations\textquotedblright{}.) Say it once more.
\end{tcolorbox}
\section[\texorpdfstring{\mr{रांगोळी} (rāṅgoḷī)}{रांगोळी (rāṅgoḷī)}]{\mr{रांगोळी} (rāṅgoḷī) — rangoli; a floor pattern}

\label{lesson:MR-C319-rangoli}



\begin{quote}
Before the new one: say the Marathi for an invitation card, then the Marathi for festive lights.
\end{quote}

\subsection*{You'll want to know: \mr{रांगोळी}}
\textbf{\mr{रांगोळी}} — \emph{rāṅgoḷī} — \textquotedblleft{}rangoli; a floor pattern\textquotedblright{}.

\begin{quote}
\mr{आईने} \mr{दारासमोर} \mr{सुंदर} \mr{रांगोळी} \mr{काढली}.
\end{quote}

\begin{quote}
\emph{āīne dārāsamor sundar rāṅgoḷī kāḍhlī.}
\end{quote}

\begin{quote}
Mother made a beautiful rangoli in front of the door.
\end{quote}

Feminine.

\mr{रांगोळी} \mr{काढणे}, to make one with coloured powder.

\begin{cousinweb}
\textbf{\mr{रांगोळी}} (\emph{rāṅgoḷī}) is the Marathi word for a pattern of coloured powder on the floor; its origin is not settled.
\end{cousinweb}

\begin{grammarlens}[title={What you've built}]
One more word for telling what happened.
\end{grammarlens}

\subsection*{Your turn}
\begin{itemize}
\raggedright
  \item \emph{Say it:} \emph{rāṅgoḷī}
  \item \emph{Say it:} \emph{rāṅgoḷī}, once more
  \item \emph{Say it:} \emph{āīne dārāsamor sundar rāṅgoḷī kāḍhlī.} — Mother made a beautiful rangoli in front of the door.
  \item \emph{Recall:} say the Marathi for an invitation card, then the Marathi for festive lights, then say \emph{rāṅgoḷī} again
\end{itemize}

\begin{tcolorbox}[breakable,colback=teal!4,colframe=teal!35!black,title={Before you move on}]
What does \mr{रांगोळी} mean? (\textquotedblleft{}Rangoli; a floor pattern\textquotedblright{}.) Say it once more.
\end{tcolorbox}
\section[\texorpdfstring{\mr{तोरण} (toraṇ)}{तोरण (toraṇ)}]{\mr{तोरण} (toraṇ) — a door garland}

\label{lesson:MR-C319-toran}



\begin{quote}
Before the new one: say the Marathi for festive lights, then the Marathi for rangoli.
\end{quote}

\subsection*{You'll want to know: \mr{तोरण}}
\textbf{\mr{तोरण}} — \emph{toraṇ} — \textquotedblleft{}a door garland\textquotedblright{}.

\begin{quote}
\mr{सणाच्या} \mr{दिवशी} \mr{दाराला} \mr{आंब्याच्या} \mr{पानांचं} \mr{तोरण} \mr{लावलं}.
\end{quote}

\begin{quote}
\emph{saṇācyā divśī dārālā āmbyācyā pānāñcaṁ toraṇ lāvlaṁ.}
\end{quote}

\begin{quote}
On the festival day a garland of mango leaves was hung on the door.
\end{quote}

Neuter.

\mr{दाराला} \mr{तोरण} \mr{बांधणे}, to hang a festoon over the door.

\begin{cousinweb}
\textbf{\mr{तोरण}} (\emph{toraṇ}) is a tatsama, taken from Sanskrit \emph{toraṇa}, 'an arched doorway, a festive gateway' (Monier-Williams).
\end{cousinweb}

\begin{grammarlens}[title={What you've built}]
One more word for telling what happened.
\end{grammarlens}

\subsection*{Your turn}
\begin{itemize}
\raggedright
  \item \emph{Say it:} \emph{toraṇ}
  \item \emph{Say it:} \emph{toraṇ}, once more
  \item \emph{Say it:} \emph{saṇācyā divśī dārālā āmbyācyā pānāñcaṁ toraṇ lāvlaṁ.} — On the festival day a garland of mango leaves was hung on the door.
  \item \emph{Recall:} say the Marathi for festive lights, then the Marathi for rangoli, then say \emph{toraṇ} again
\end{itemize}

\begin{tcolorbox}[breakable,colback=teal!4,colframe=teal!35!black,title={Before you move on}]
What does \mr{तोरण} mean? (\textquotedblleft{}A door garland\textquotedblright{}.) Say it once more.
\end{tcolorbox}
\section[\texorpdfstring{\mr{फटाका} (phaṭākā)}{फटाका (phaṭākā)}]{\mr{फटाका} (phaṭākā) — a firecracker}

\label{lesson:MR-C319-phataka}



\begin{quote}
Before the new one: say the Marathi for rangoli, then the Marathi for a door garland.
\end{quote}

\subsection*{You'll want to know: \mr{फटाका}}
\textbf{\mr{फटाका}} — \emph{phaṭākā} — \textquotedblleft{}a firecracker\textquotedblright{}.

\begin{quote}
\mr{रात्री} \mr{मुलांनी} \mr{खूप} \mr{फटाके} \mr{फोडले}.
\end{quote}

\begin{quote}
\emph{rātrī mulānnī khūp phaṭāke phoḍle.}
\end{quote}

\begin{quote}
At night the children set off lots of firecrackers.
\end{quote}

Masculine. Plural: \mr{फटाके}, \emph{phaṭāke}.

Plural \mr{फटाके}; \mr{फटाके} \mr{फोडणे}, to set off crackers.

\begin{cousinweb}
\textbf{\mr{फटाका}} (\emph{phaṭākā}) is a Marathi word that seems built on the sound \emph{phaṭ}, a bang; its origin is not settled.
\end{cousinweb}

\begin{grammarlens}[title={What you've built}]
One more word for telling what happened.
\end{grammarlens}

\subsection*{Your turn}
\begin{itemize}
\raggedright
  \item \emph{Say it:} \emph{phaṭākā}
  \item \emph{Say it:} \emph{phaṭākā}, once more
  \item \emph{Say it:} \emph{rātrī mulānnī khūp phaṭāke phoḍle.} — At night the children set off lots of firecrackers.
  \item \emph{Recall:} say the Marathi for rangoli, then the Marathi for a door garland, then say \emph{phaṭākā} again
\end{itemize}

\begin{tcolorbox}[breakable,colback=teal!4,colframe=teal!35!black,title={Before you move on}]
What does \mr{फटाका} mean? (\textquotedblleft{}A firecracker\textquotedblright{}.) Say it once more.
\end{tcolorbox}
\section[\texorpdfstring{\mr{आतषबाजी} (ātaṣbājī)}{आतषबाजी (ātaṣbājī)}]{\mr{आतषबाजी} (ātaṣbājī) — fireworks}

\label{lesson:MR-C319-atasbaji}



\begin{quote}
Before the new one: say the Marathi for a door garland, then the Marathi for a firecracker.
\end{quote}

\subsection*{You'll want to know: \mr{आतषबाजी}}
\textbf{\mr{आतषबाजी}} — \emph{ātaṣbājī} — \textquotedblleft{}fireworks\textquotedblright{}.

\begin{quote}
\mr{जत्रेच्या} \mr{शेवटी} \mr{आकाशात} \mr{आतषबाजी} \mr{झाली}.
\end{quote}

\begin{quote}
\emph{jatrecyā śevṭī ākāśāt ātaṣbājī jhālī.}
\end{quote}

\begin{quote}
At the end of the fair there were fireworks in the sky.
\end{quote}

Feminine.

\mr{आतषबाजी} \mr{करणे}, to put on fireworks.

\begin{cousinweb}
\textbf{\mr{आतषबाजी}} (\emph{ātaṣbājī}) is a loan from Persian \emph{ātaš-bāzī}, built from \emph{ātaš}, 'fire', and \emph{bāzī}, 'play'.
\end{cousinweb}

\begin{grammarlens}[title={What you've built}]
One more word for telling what happened.
\end{grammarlens}

\subsection*{Your turn}
\begin{itemize}
\raggedright
  \item \emph{Say it:} \emph{ātaṣbājī}
  \item \emph{Say it:} \emph{ātaṣbājī}, once more
  \item \emph{Say it:} \emph{jatrecyā śevṭī ākāśāt ātaṣbājī jhālī.} — At the end of the fair there were fireworks in the sky.
  \item \emph{Recall:} say the Marathi for a door garland, then the Marathi for a firecracker, then say \emph{ātaṣbājī} again
\end{itemize}

\begin{tcolorbox}[breakable,colback=teal!4,colframe=teal!35!black,title={Before you move on}]
What does \mr{आतषबाजी} mean? (\textquotedblleft{}Fireworks\textquotedblright{}.) Say it once more.
\end{tcolorbox}
\section[(dialogue)]{Looking back: from On the Air to Lights and Colours}

\label{lesson:MR-R319-review-312-319}



\begin{quote}
Say the Marathi for radio and television.
\end{quote}

\subsection*{Your turn}
Say these aloud:

\begin{itemize}
\raggedright
  \item radio, television, a channel, a newsreader, to broadcast
  \item a rumour, a sensation, an announcement, prominent, to proclaim
  \item a reader, a viewer, a listener, a title, a front page
  \item an eyewitness, a documentary, a magazine issue, to register, to wander
  \item a wedding, an engagement, henna, a wedding procession, a marquee
  \item a wedding gift, a blessing, blessing rice, a shehnai, a rhyming couplet
  \item birth, a naming ceremony, a birthday, a thread ceremony, an invitation card
  \item festive lights, rangoli, a door garland, a firecracker, fireworks
\end{itemize}

\begin{tcolorbox}[breakable,colback=teal!4,colframe=teal!35!black,title={Before you move on}]
Radio? (\textbf{\mr{आकाशवाणी}}, \emph{ākāśvāṇī}.) Television? (\textbf{\mr{दूरदर्शन}}, \emph{dūrdarśan}.) A channel? (\textbf{\mr{वाहिनी}}, \emph{vāhinī}.) A newsreader? (\textbf{\mr{निवेदक}}, \emph{nivedak}.) To broadcast? (\textbf{\mr{प्रसारित} \mr{करणे}}, \emph{prasārit karṇe}.) A rumour? (\textbf{\mr{अफवा}}, \emph{aphvā}.) A sensation? (\textbf{\mr{खळबळ}}, \emph{khaḷbaḷ}.) An announcement? (\textbf{\mr{घोषणा}}, \emph{ghoṣṇā}.) Prominent? (\textbf{\mr{ठळक}}, \emph{ṭhaḷak}.) To proclaim? (\textbf{\mr{पुकारणे}}, \emph{pukārṇe}.) A reader? (\textbf{\mr{वाचक}}, \emph{vācak}.) A viewer? (\textbf{\mr{प्रेक्षक}}, \emph{prekṣak}.) A listener? (\textbf{\mr{श्रोता}}, \emph{śrotā}.) A title? (\textbf{\mr{शीर्षक}}, \emph{śīrṣak}.) A front page? (\textbf{\mr{मुखपृष्ठ}}, \emph{mukhpṛṣṭha}.) An eyewitness? (\textbf{\mr{प्रत्यक्षदर्शी}}, \emph{pratyakṣadarśī}.) A documentary? (\textbf{\mr{माहितीपट}}, \emph{māhitīpaṭ}.) A magazine issue? (\textbf{\mr{अंक}}, \emph{aṅk}.) To register? (\textbf{\mr{नोंदवणे}}, \emph{nondavṇe}.) To wander? (\textbf{\mr{भटकणे}}, \emph{bhaṭakṇe}.) A wedding? (\textbf{\mr{लग्न}}, \emph{lagna}.) An engagement? (\textbf{\mr{साखरपुडा}}, \emph{sākharpuḍā}.) Henna? (\textbf{\mr{मेंदी}}, \emph{mendī}.) A wedding procession? (\textbf{\mr{वरात}}, \emph{varāt}.) A marquee? (\textbf{\mr{मंडप}}, \emph{maṇḍap}.) A wedding gift? (\textbf{\mr{आहेर}}, \emph{āher}.) A blessing? (\textbf{\mr{आशीर्वाद}}, \emph{āśīrvād}.) Blessing rice? (\textbf{\mr{अक्षता}}, \emph{akṣatā}.) A shehnai? (\textbf{\mr{सनई}}, \emph{sanaī}.) A rhyming couplet? (\textbf{\mr{उखाणा}}, \emph{ukhāṇā}.) Birth? (\textbf{\mr{जन्म}}, \emph{janma}.) A naming ceremony? (\textbf{\mr{बारसे}}, \emph{bārse}.) A birthday? (\textbf{\mr{वाढदिवस}}, \emph{vāḍhdivas}.) A thread ceremony? (\textbf{\mr{मुंज}}, \emph{muñj}.) An invitation card? (\textbf{\mr{पत्रिका}}, \emph{patrikā}.) Festive lights? (\textbf{\mr{रोषणाई}}, \emph{roṣṇāī}.) Rangoli? (\textbf{\mr{रांगोळी}}, \emph{rāṅgoḷī}.) A door garland? (\textbf{\mr{तोरण}}, \emph{toraṇ}.) A firecracker? (\textbf{\mr{फटाका}}, \emph{phaṭākā}.) Fireworks? (\textbf{\mr{आतषबाजी}}, \emph{ātaṣbājī}.)
\end{tcolorbox}
